\section{液-固界面吸附}


\begin{frame}
    \frametitle{溶液中吸附量的测定}

    固体在溶液中的吸附较为复杂，因为吸附剂既要吸附溶质，又要吸附溶剂，溶质和溶剂之间还有相互作用。
    
    溶液中的吸附虽然很复杂，但吸附量的测定却比较简单。通常是在定温条件下，将一定量的固体吸附剂与一定量已知浓度的溶液混合，达到吸附平衡后分析溶液的成分。从吸附前后溶液浓度的改变可求出固体对溶质的吸附量。
    \[
        \varGamma = \frac{x}{m} = \frac{V(c_1 - c_2)}{m}
    \]
    式中：$m$为吸附剂的质量；$c_1$和$c_2$分别为溶液起始和平衡的浓度；$V$为溶液的体积，$x$为被$m$g吸附剂所吸附溶质的物质的量。$\varGamma$实际上是固体吸附溶质、溶剂的总结果，称为表观或相对吸附量。 
    
    如果溶液很稀，固体吸附溶剂所造成的浓度变化可以忽略，求得的吸附量可近似看作溶质的吸附量。

\end{frame}


\begin{frame}
    \frametitle{稀溶液中溶质分子的吸附}

    \begin{itemize}
        \item[1.] \cbf{吸附定温线}
    \end{itemize}
    对于不同的体系，吸附定温线的形状不同，有些液-固体系可以使用某些气-固吸附的定温式。
    
    用弗兰德里希定温式表示稀溶液中吸附量与平衡浓度的关系
    \[
        \varGamma = \frac{x}{m} = kc^{\frac{1}{n}}
    \]
    有些稀溶液体系可以用朗缪尔方程来描述其吸附定温线
    \[
        \varGamma = \frac{x}{m} = \frac{\varGamma_\mathrm{m} bc}{1 + bc}
    \]
    如果测出$\varGamma_\mathrm{m}$值和溶质分子截面积$a$，则可以估算出吸附剂的比表积$S_0$
    \[
        S_0 = \varGamma_\mathrm{m} La
    \]

\end{frame}


\begin{frame}
    \frametitle{稀溶液中溶质分子的吸附}
    \begin{itemize}
        \item[2.] \cbf{影响吸附的因素}
    \end{itemize}
    由于溶液中的各组分都能被固体表面吸附，组分之间又存在着相互作用，因此影响吸附的因素较复杂。在长期的实践中，总结出了一些经验规律：
    \begin{itemize}
        \item \bf{极性的影响}
        
        一般说来，极性吸附剂在非极性溶剂中优先吸附极性强的溶质，而非极性的吸附剂在极性溶剂中优先吸附非极性强的溶质。
        \item \bf{溶质溶解度的影响}
        
        溶质溶解度越小，说明溶质与溶剂之间的相互作用越弱，溶质越容易从溶液中析出，较易被吸附。
        \item \bf{温度的影响}
        
        一般情况下，温度升高，吸附量下降。
    \end{itemize}

\end{frame}


\begin{frame}
    \frametitle{电解质溶液中离子的吸附}

    固体在电解质溶液中对离子的吸附通常有两种：
    \begin{itemize}
        \item \bf{离子选择性吸附}
        
        固体在电解质溶液中，可以优先地吸附某种电荷的离子（正离子或负离子），而使固体带有电荷（正或负），这种现象称为离子的选择性吸附。
        \item \bf{离子交换吸附}
        
        离子交换吸附不是固体直接从溶液中吸附离子，而是固体吸附剂本身的组成离子与溶液中的同号离子发生交换，即固体吸附剂在溶液中吸附了某种离子的同时，将另一种相同符号电荷的离子释放到溶液中。
    \end{itemize}

\end{frame}